% AI in Genomics -- Project Writeup
\documentclass[11pt]{article}
\usepackage[utf8]{inputenc}
\usepackage{geometry}
\usepackage{hyperref}
\usepackage{graphicx}
\usepackage{booktabs}
\usepackage{enumitem}
\usepackage{caption}
\geometry{margin=1in}
\graphicspath{{../figures/}{../}{./}}
\emergencystretch=3em

\title{Environmental Conditioning of Protein Sequence Generation\\
       from GOLD, UniProt, NASA POWER, and WorldClim}
\author{AI in Genomics Project}
\date{}

\begin{document}
\maketitle

\begin{abstract}
This project asks whether environmental measurements attached to biological
sampling locations can be used to condition a protein sequence generator. The
implemented code builds a dataset by linking GOLD organism metadata to reviewed
UniProt protein sequences, enriching sample coordinates with NASA POWER daily
weather and WorldClim monthly climate rasters, cleaning invalid numerical
sentinels, and training a recurrent neural network that generates amino-acid
sequences from environmental inputs. The current model should be treated as a
proof of concept: it demonstrates an end-to-end reproducible pipeline and an
interactive Streamlit interface, but it does not yet prove that generated
sequences are biologically functional or causally adapted to the environment.
\end{abstract}

\section{Background}
Proteins are chains of amino acids that fold into molecular machines. Their
sequences are constrained by evolution, organism history, and environment. For
example, proteins from organisms living in hot or dry habitats may be under
different selective pressures than proteins from temperate habitats. Our goal is
not to model all of evolution, but to test a narrower machine-learning question:
can a model learn a relationship between coarse environmental variables and the
distribution of protein sequences observed in organisms sampled from those
environments?

The project combines four public resources. GOLD provides organism-level sample
metadata, including taxonomy IDs and geographic coordinates when available
\cite{gold}. UniProt provides reviewed protein sequences for each taxonomy ID
\cite{uniprot}. NASA POWER provides daily surface temperature and radiation for
the sample location/date \cite{nasa}. WorldClim provides long-term monthly
climate rasters, including precipitation, solar radiation, average temperature,
and vapour pressure \cite{worldclim}. SoilGrids querying is implemented in the
code, but the current merged training table does not include the returned pH or
soil-carbon fields as model inputs \cite{soilgrids}.

\section{Motivation}
Most protein generation models are conditioned on sequence context, structure,
or textual prompts. We explored a different conditioning signal: the physical
environment associated with the organism from which a protein was sampled. If
such a signal is useful, it could eventually help generate candidate protein
sequences for unusual environmental settings. The current project tests the
infrastructure and modeling feasibility of this idea rather than claiming a
validated biological design tool.

\section{Data}
\subsection{Source Data and Linking}
The pipeline begins with the \texttt{Organism} sheet in \texttt{goldData.xlsx}.
The loader caches this sheet as \texttt{goldData\_Organism.parquet}. In the
local cached copy, GOLD contains 530,160 organism rows; 64,665 rows have both
latitude and longitude. The code sorts these coordinate-bearing rows and queries
UniProt using each row's NCBI taxonomy ID.

For each taxonomy ID, \texttt{src/main.py} calls the UniProt REST search API for
up to five reviewed protein entries. The current query is
\texttt{taxonomy\_id:<id> AND reviewed:true}; it retrieves accession, entry
name, organism name, and sequence. These protein rows are then joined back to
GOLD by taxonomy ID and enriched by geographic location.

\subsection{Environmental Covariates}
Each unique latitude, longitude, and collection-date triple is enriched with:
\begin{itemize}[nosep]
    \item NASA POWER daily temperature at 2 meters and shortwave radiation.
    These are stored as \texttt{NASA\_Temp\_C} and \texttt{NASA\_Radiation}.
    If GOLD lacks a parseable collection date, the code uses 2015-01-01 as a
    fallback date.
    \item WorldClim monthly precipitation (\texttt{prec\_01}--\texttt{prec\_12}),
    solar radiation (\texttt{srad\_01}--\texttt{srad\_12}), average
    temperature (\texttt{tavg\_01}--\texttt{tavg\_12}), and vapour pressure
    (\texttt{vapr\_01}--\texttt{vapr\_12}) sampled from local GeoTIFF rasters.
    \item SoilGrids pH and soil organic carbon are fetchable through
    \texttt{src/soilgrids\_data.py}, but those fields are currently commented
    out during dataset assembly and are not used by the trained model.
\end{itemize}

\subsection{Cleaning and Dataset Size}
\texttt{src/clean\_data.py} converts the merged parquet to CSV and removes rows
containing common no-data or overflow sentinels in any numeric column, including
\texttt{-9999}, \texttt{-32768}, \texttt{32767}, \texttt{65535}, or values
larger than \texttt{1e30} in magnitude. Table~\ref{tab:data} summarizes the
largest cleaned dataset currently present in the repository.

\begin{table}[htbp]
\centering
\caption{Current local dataset used by the implemented training pipeline. The
``raw merged'' row refers to \texttt{combined\_uniprot\_and\_gold\_environmental\_data\_10K.csv};
the ``cleaned'' row refers to \texttt{cleaned\_environmental\_data\_10K.csv}.}
\label{tab:data}
\begin{tabular}{lrrrr}
\toprule
Dataset & Rows & Unique taxonomy IDs & Unique UniProt accessions & Median sequence length \\
\midrule
Raw merged & 10,113 & 830 & 8,888 & 381 aa \\
Cleaned & 8,631 & 720 & 7,809 & 381 aa \\
\bottomrule
\end{tabular}
\end{table}

In the cleaned 10K dataset, 8,376 rows have NASA temperature, and 8,296 rows
have NASA radiation. Protein sequence lengths range from 6 to 11,872 amino
acids, with a mean length of 450 amino acids. During training, sequences are
truncated to 198 amino acids plus start and end tokens, so very long proteins
are represented only by their N-terminal prefix.

\section{Methods}
\subsection{Feature Engineering}
The model uses three environmental input features:
\begin{enumerate}[nosep]
    \item \texttt{NASA\_Temp\_C}: daily NASA POWER temperature for the sample
    location/date.
    \item \texttt{Total\_Annual\_Prec}: sum of the 12 WorldClim monthly
    precipitation values.
    \item \texttt{Mean\_Annual\_Solar\_Rad}: mean of the 12 WorldClim monthly
    solar-radiation values.
\end{enumerate}
Although \texttt{src/train.py} also computes \texttt{Mean\_Annual\_Temp}, that
field is not included in the current model input. The three selected numeric
features are standardized with scikit-learn's \texttt{StandardScaler}. The
fitted scaler is saved as \texttt{checkpoints/environmental\_scaler.pkl} and is
reused by the Streamlit app.

\subsection{Sequence Encoding}
The output sequence is a character-level amino-acid sequence. The vocabulary has
23 tokens: the 20 standard amino acids plus \texttt{<PAD>}, \texttt{<SOS>}, and
\texttt{<EOS>}. Each training sequence is truncated to a maximum length of 200
tokens after adding start and end markers. Unknown or nonstandard amino-acid
characters are currently mapped to \texttt{<PAD>}; this is simple, but it means
rare sequence characters are treated like padding rather than modeled
explicitly.

\subsection{Model Architecture}
The implemented model in \texttt{src/model.py} is a conditional LSTM protein
generator. The standardized three-dimensional environmental vector is passed
through two learned linear layers. One initializes the LSTM hidden state and the
other initializes the LSTM cell state. The amino-acid tokens are embedded into
64-dimensional vectors, processed by a single-layer LSTM with hidden size 256,
and projected to the 23-token vocabulary at each position. During training, the
model learns next-token prediction: given the environment and all previous
tokens, it predicts the next amino acid.

For comparison, \texttt{src/train\_baseline.py} trains a simpler conditional
RNN baseline with the same environmental inputs, token vocabulary, embedding
size, hidden size, optimizer, and train/validation split. The baseline receives
the environmental vector only as the initial recurrent hidden state, while the
LSTM model uses both a hidden and cell state.

\subsection{Training and Evaluation Protocol}
\texttt{src/train.py} drops rows missing sequence, NASA temperature, annual
precipitation, or annual solar radiation. It then splits rows into 90\% training
and 10\% validation using \texttt{random\_state=42}. Training uses Adam with
learning rate 0.001 for 100 epochs, batch size 32, and cross-entropy loss over
next-token predictions. Padding tokens are ignored in the loss. The best model
by validation loss is saved to the requested checkpoint path.

\texttt{src/evaluation.py} evaluates the saved LSTM and RNN checkpoints on the
same validation split. It reports cross-entropy loss, perplexity, next-token
accuracy excluding padding, and GRAVY scores for generated hot and cold
environment examples. This split prevents exact row overlap between training and
validation, but it is not a stringent biological generalization test because
related organisms or protein families can appear on both sides.

\section{Results}
\subsection{End-to-End Pipeline}
The main result is a working end-to-end pipeline. \texttt{src/main.py} builds a
merged protein/environment table, \texttt{src/clean\_data.py} removes corrupted
numeric rows, \texttt{src/train.py} trains and checkpoints the conditional LSTM,
and \texttt{src/app.py} loads the checkpoint and scaler into an interactive
Streamlit interface. The repository currently contains trained checkpoints for
the 5K and 10K data runs, plus a fitted environmental scaler.

\subsection{Sequence Space Visualization}
\texttt{src/visualization\_pca.py} and \texttt{src/visualization\_tsne.py}
provide exploratory checks of whether the sequence dataset contains visible
environmental structure. Both scripts convert each protein sequence into 3-mer
TF-IDF features and color the resulting points by environmental variables. PCA
gives a linear low-dimensional summary, while t-SNE first reduces the TF-IDF
matrix to 50 dimensions with truncated SVD and then emphasizes local sequence
neighborhoods in two dimensions. These plots are not model-evaluation metrics,
but they help reveal climate gradients or data artifacts in sequence feature
space (Figures~\ref{fig:sequence-space-pca} and \ref{fig:sequence-space-tsne}).

\begin{figure}[ht]
\centering
\centering
\includegraphics[width=\linewidth]{environmental_sequence_pca_5plots.png}
\caption{Exploratory sequence-space PCA visualization. PCA gives a linear
two-dimensional summary of 3-mer TF-IDF protein features, so broad gradients
reflect large-scale differences in amino-acid 3-mer composition. In the
plot, each point is one protein sequence, and the five panels color the same
points by NASA surface temperature, WorldClim annual temperature, annual
precipitation, annual solar radiation, and NASA radiation.}
\label{fig:sequence-space-pca}
\end{figure}

\begin{figure}[ht]
\centering
\includegraphics[width=\linewidth]{environmental_sequence_tsne_5plots.png}
\caption{Exploratory sequence-space TSNE visualizations. t-SNE
uses the summary of 3-mer TF-IDF features but emphasizes local neighborhoods, so compact
clusters indicate sequences with similar short amino-acid patterns. In the
plot, each point is one protein sequence, and the five panels color the same
points by NASA surface temperature, WorldClim annual temperature, annual
precipitation, annual solar radiation, and NASA radiation.}
\label{fig:sequence-space-tsne}
\end{figure}

The TSNE plots show some visible clustering and climate gradients, but the PCA plots do not show much structure. This suggests that local sequence neighborhoods may be more environmentally structured than global linear projections of sequence space. However, these visualizations are not rigorous evaluations of model performance; they are only intended to reveal broad patterns in the training data.

\subsection{Baseline Comparison}
The LSTM outperformed the simpler RNN baseline on the validation split
(Table~\ref{tab:baseline}). It achieved lower cross-entropy loss, lower
perplexity, and higher next-token accuracy. For generated sequences, both models
produced length-100 examples under the tested hot and cold environments. GRAVY
is the grand average of hydropathy: positive values indicate more hydrophobic
sequences, while negative values indicate more hydrophilic sequences. The LSTM
generated a hydrophobic hot-environment sequence and a hydrophilic
cold-environment sequence, while the RNN baseline generated strongly
hydrophobic sequences in both cases.

\begin{table}[htbp]
\centering
\caption{Validation and generated-sequence metrics from
\texttt{results/baseline\_vs\_lstm.txt}. Statistical metrics were computed on
the 10\% validation split; GRAVY metrics were computed from one generated
sequence per model for each test environment.}
\label{tab:baseline}
\begin{tabular}{lrrr}
\toprule
Model & Loss & Perplexity & Next-token accuracy \\
\midrule
LSTM & 2.4350 & 11.42 & 25.33\% \\
RNN baseline & 2.6927 & 14.77 & 18.48\% \\
\bottomrule
\end{tabular}

\vspace{0.75em}
\begin{tabular}{llrr}
\toprule
Environment & Model & Generated length & GRAVY \\
\midrule
Hot, 85$^\circ$C & LSTM & 100 & 0.434 \\
Hot, 85$^\circ$C & RNN baseline & 100 & 1.042 \\
Cold, -15$^\circ$C & LSTM & 100 & -0.419 \\
Cold, -15$^\circ$C & RNN baseline & 100 & 1.979 \\
\bottomrule
\end{tabular}
\end{table}

\subsection{Generation Interface}
The trained model can generate a sequence from user-specified temperature,
annual precipitation, and solar-radiation values. In command-line inference,
\texttt{src/train.py} tests three example environments: an Earth-like temperate
case, a hot arid high-radiation case, and a frozen low-radiation case.
\texttt{src/app.py} exposes the same generation function in Streamlit with
sliders and preset environments. Generation is greedy: at each step, the model
chooses the most likely next token until it emits \texttt{<EOS>} or reaches the
maximum generation length.

\section{Discussion and Limitations}
The project demonstrates that a heterogeneous genomics-and-environment dataset
can be assembled and used to condition a sequence generator. The codebase now
has modular data retrieval, cleaning, training, checkpointing, visualization,
and UI components. This is the strongest outcome of the current work: the
pipeline is reproducible and can be iterated on.

The biological interpretation is still limited. First, the input features are
coarse descriptors of sampling location, not direct measurements of the
microenvironment where every organism grew. Second, proteins are strongly shaped
by function and ancestry; a random row split can overestimate performance if
closely related organisms or similar protein families appear in both training
and validation. Third, the model generates short prefixes of proteins because
training truncates sequences to 200 tokens. Finally, the current evaluation
contains next-token and simple hydropathy metrics, but not edit distance,
protein-family holdout tests, or functional validation. Those metrics should be
added before making strong claims about model quality.

The next most important improvements are therefore methodological rather than
architectural: save training curves, split data by taxonomy ID or protein
family, add sequence-quality filters for nonstandard amino acids, compare
against an unconditional sequence model, and evaluate generated sequences with
additional protein-level summaries such as amino-acid composition and predicted
stability. SoilGrids features are also already implemented and could be added to
the model input once their coverage and units are checked.

\section{Reproducibility}
The main commands are:
\begin{verbatim}
python src/main.py

python src/clean_data.py \
  --input-parquet combined_uniprot_and_gold_environmental_data_10K.parquet \
  --output-csv combined_uniprot_and_gold_environmental_data_10K.csv \
  --cleaned-output-csv cleaned_environmental_data_10K.csv

python src/train.py \
  --data cleaned_environmental_data_10K.csv \
  --checkpoint checkpoints/best_alien_protein_model_10K.pt

streamlit run src/app.py -- \
  --checkpoint checkpoints/best_alien_protein_model_10K.pt \
  --scaler checkpoints/environmental_scaler.pkl
\end{verbatim}
The WorldClim GeoTIFFs are expected under \texttt{data/worldclim/}. The main
scripts are organized under \texttt{src/}; checkpoints and the scaler are stored
under \texttt{checkpoints/}; exploratory figures are stored under
\texttt{figures/}; evaluation summaries are stored under \texttt{results/}.
External API calls require internet access and may change if source services
update their schemas or availability.

\begin{thebibliography}{9}
\bibitem{gold}
GOLD: Genomes Online Database. \url{https://gold.jgi.doe.gov/}

\bibitem{uniprot}
The UniProt Consortium. UniProt: the Universal Protein Knowledgebase.
\url{https://www.uniprot.org/}

\bibitem{nasa}
NASA POWER Project. Prediction Of Worldwide Energy Resources.
\url{https://power.larc.nasa.gov/}

\bibitem{worldclim}
WorldClim version 2.1 climate data. \url{https://www.worldclim.org/}

\bibitem{soilgrids}
ISRIC SoilGrids. \url{https://soilgrids.org/}
\end{thebibliography}

\end{document}
